% !PRINTED 446
\weberpage{446}
%<<CONT
befriedigt für
\[
u = (k), \quad v = (lk) \quad \text{und} \quad v = (l^{-1}k).
\]
Nach dem, was am Schlusse des § 69 bewiesen ist [Formel
(32)], folgt aber hieraus die Kongruenz
\weq{16}{[(k)^p - (lk)] [(lk)^p - (k)] \equiv 0 \quad (\mathrm{mod}\ p).}
Wenn nun $\mathfrak P$ ein in $p$ aufgehendes Primideal ist, so muß
einer der beiden Faktoren auf der linken Seite von (17) % [sic] printed (17); only (16) has two factors on its left side
durch $\mathfrak P$
teilbar sein. Welche der beiden hiernach möglichen Annahmen
wir weiter verfolgen, ist gleichgültig, da die eine in die andere
übergeht, wenn $k$ mit $l^{-1}k$ und $l$ mit $l^{-1}$ vertauscht wird. Sei
also
\weq{17}{(k)^p \equiv (lk) \quad (\mathrm{mod}\ \mathfrak P).}
Wenn die Kongruenz (17) für irgend einen der konjugierten
Werte $(k)$ befriedigt ist, so gilt sie auch für jeden anderen.
Denn es ist nach § 121:
\[
(k') = f(k, \sqrt{D}), \quad (lk') = f(lk, \sqrt{D}),
\]
also
\[
(k')^p \equiv f(k^p, \sqrt{D}) \equiv f(lk, \sqrt{D}) \equiv (lk') \quad (\mathrm{mod}\ \mathfrak P).
\]
Es folgt also aus (14), % [sic] printed (14); the substitution described applies to (17)
wenn man $k$ durch $l^{-1}k$ ersetzt:
\weq{18}{(l^{-1}k)^p \equiv (k) \quad (\mathrm{mod}\ \mathfrak P),}
d. h. je nachdem für $l$ die eine oder die andere der beiden zu $p$
gehörigen Klassen $l$ gesetzt wird, ist der eine oder der andere
Faktor von (16) durch $\mathfrak P$ teilbar.
Durch wiederholte Anwendung von (17) ergibt sich für jeden
beliebigen positiven Exponenten $\pi$
\weq{19}{(k)^{p^\pi} \equiv (l^\pi k) \quad (\mathrm{mod}\ \mathfrak P)}
(Schluß von $\pi$ auf $\pi + 1$).
Wenn nun $l$ zum Exponenten $\varepsilon$ gehört, oder $\varepsilon$ der Index von
$p$ ist, so ist nach (19):
\weq{20}{(k)^{p^\varepsilon} \equiv (k) \quad (\mathrm{mod}\ \mathfrak P),}
und $p^\varepsilon$ ist die \gesperrt{niedrigste} Potenz von $p$, die dieser Bedingung
genügt. Denn wenn noch eine niedrigere Potenz von $p$ die Kongruenz (19) erfüllt, % [sic] printed references (19) then (20); their logical roles appear interchanged
so folgt aus (20), daß zwei \gesperrt{verschiedene}
Klasseninvarianten $(k)$, $(k')$ nach dem Modul $\mathfrak P$ kongruent sind.
Es wäre also ihre Differenz $(k) - (k')$ durch $\mathfrak P$ und mithin die
Diskriminante von $H_m$ durch $p$ teilbar, gegen die Voraussetzung.
